% GENERATED FILE. Edit canonical lessons, then run npm run generate:books.
% canonical-source-hash: fnv1a64:b8ed939075f5a8f3
% canonical-lessons: ZH-C381-fenlei, ZH-C381-quanjia, ZH-W381-quan-whole, ZH-C381-huishou, ZH-C381-quanguo, ZH-C381-anquangan

\chapter{The Whole Family}
\label{ch:zh-the-whole-family}
\hlchaptermodality{381}

\begin{chapteropening}
\textbf{By the end of this chapter:} \emph{I can talk about sorting rubbish, the whole family, recycling, the whole country and feeling safe.}

Complete the last lesson of chapter 381: I can talk about sorting rubbish, the whole family, recycling, the whole country and feeling safe.
\end{chapteropening}

\section[fēnlèi]{fēnlèi — to sort into kinds}

\label{lesson:ZH-C381-fenlei}



\begin{quote}
Before the new one: say the Mandarin for our planet, then the Mandarin for eco-friendly (good for the environment).
\end{quote}

\subsection*{You'll want to know: fēnlèi}
\emph{fēnlèi} — \textquotedblleft{}to sort into kinds\textquotedblright{}.

\begin{quote}
\emph{lājī yīnggāi fēnlèi} — Rubbish should be sorted.
\end{quote}

Its characters come later; for now it is a word for the ear.

\begin{cousinweb}
\emph{fēnlèi} is built from \emph{fēn} 'to divide' and \emph{lèi} 'a kind', both inherited from Middle Chinese, \emph{pjun} and \emph{lwijH} (Guangyun, in Baxter's transcription), which glosses the first 'to separate' and the second 'a sort'.
\end{cousinweb}

\begin{grammarlens}[title={What you've built}]
One more word for giving reasons.
\end{grammarlens}

\subsection*{Your turn}
\begin{itemize}
\raggedright
  \item \emph{Say it:} \emph{fēnlèi}
  \item \emph{Say it:} \emph{fēnlèi}, once more
  \item \emph{Say it:} \emph{lājī yīnggāi fēnlèi} — Rubbish should be sorted.
  \item \emph{Recall:} say the Mandarin for our planet, then the Mandarin for eco-friendly (good for the environment), then say \emph{fēnlèi} again
\end{itemize}

\begin{tcolorbox}[breakable,colback=teal!4,colframe=teal!35!black,title={Before you move on}]
What does fēnlèi mean? (\textquotedblleft{}To sort into kinds\textquotedblright{}.) Say it once more.
\end{tcolorbox}
\section[\texorpdfstring{\zh{全家} (quánjiā)}{全家 (quánjiā)}]{\zh{全家} (quánjiā) — the whole family}

\label{lesson:ZH-C381-quanjia}



\begin{quote}
Before the new one: say the Mandarin for eco-friendly (good for the environment), then the Mandarin for to sort into kinds.
\end{quote}

\subsection*{You'll want to know: \zh{全家}}
\textbf{\zh{全家}} — \emph{quánjiā} — \textquotedblleft{}the whole family\textquotedblright{}.

\begin{quote}
\textbf{\zh{今天我们全家一起打扫房子}} \emph{jīntiān wǒmen quánjiā yìqǐ dǎsǎo fángzi} — Today the whole family is cleaning the house together.
\end{quote}

The new character, in pieces: \zh{全} — six strokes: a wide roof like \zh{人}, then a bar, a second bar, a vertical down from the first through the second, and a long bar at the foot. The next lesson puts it in your hand.

\begin{cousinweb}
\textbf{\zh{全家}} (\emph{quánjiā}) is built from \zh{全} 'whole' and \zh{家} 'a home', both inherited from Middle Chinese, \emph{dzjwen} and \emph{kae} (Guangyun, in Baxter's transcription), which glosses the first 'complete'.
\end{cousinweb}

\begin{grammarlens}[title={What you've built}]
One more word for giving reasons.
\end{grammarlens}

\subsection*{Your turn}
\begin{itemize}
\raggedright
  \item \emph{Say it:} \emph{quánjiā}
  \item \emph{Say it:} \emph{quánjiā}, once more
  \item \emph{Say it:} \emph{jīntiān wǒmen quánjiā yìqǐ dǎsǎo fángzi} — Today the whole family is cleaning the house together.
  \item \emph{Recall:} say the Mandarin for eco-friendly (good for the environment), then the Mandarin for to sort into kinds, then say \emph{quánjiā} again
\end{itemize}

\begin{tcolorbox}[breakable,colback=teal!4,colframe=teal!35!black,title={Before you move on}]
What does \zh{全家} mean? (\textquotedblleft{}The whole family\textquotedblright{}.) Say it once more.
\end{tcolorbox}
\section[\texorpdfstring{\zh{全} (quán)}{全 (quán)}]{\zh{全} — a wide roof, then three bars and a vertical}

\label{lesson:ZH-W381-quan-whole}



\begin{quote}
The word first, then the character inside it.

\begin{itemize}
\raggedright
  \item \emph{Recall:} say \emph{quánjiā} — \textbf{R1}, one lesson back, before it decays
\end{itemize}

\begin{quote}
\zh{人} \zh{安}
\end{quote}

\emph{Read it:} each one aloud, with nothing said first — \emph{rén}, \emph{ān}
\end{quote}

\subsection*{Script — a wide roof, then three bars and a vertical}
\hlblockfigure{\detokenize{figures/ZH-W381-quan-whole-filmstrip.pdf}}{How \zh{全} is written, stroke by stroke}

\begin{quote}
\zh{全}
\end{quote}

\textbf{Two pieces, top before bottom.}

\textbf{The roof, two strokes,} like \zh{人} opened wide: a long stroke falling to the left, then a long stroke falling to the right from near its top.

\textbf{Under it, four strokes.} A bar. A second bar under it. A vertical down from the first bar through the second. Then a long bar across the foot.

2 + 4 = six.

\subsection*{Your turn}
\begin{itemize}
\raggedright
  \item \emph{Write it:} \zh{全}, all six, with the model in front of you
  \item \emph{Cover:} write it again from memory, and count the pieces, not the strokes
\end{itemize}

\begin{tcolorbox}[breakable,colback=teal!4,colframe=teal!35!black,title={Before you move on}]
How many strokes in \zh{全}? (\textbf{Six}.) What is the roof like? (\textbf{\zh{人}, opened wide}.) Which stroke comes after the two bars? (\textbf{The vertical}.)

Source: \href{https://raw.githubusercontent.com/chanind/hanzi-writer-data/68d10a4b21150cae5e1ebbd223eed289cf32d90c/data/\%E5\%85\%A8.json}{Hanzi Writer Data \zh{全}, pinned stroke paths and medians}.
\end{tcolorbox}
\section[huíshōu]{huíshōu — to recycle}

\label{lesson:ZH-C381-huishou}



\begin{quote}
Before the new one: say the Mandarin for to sort into kinds, then the Mandarin for the whole family.
\end{quote}

\subsection*{You'll want to know: huíshōu}
\emph{huíshōu} — \textquotedblleft{}to recycle\textquotedblright{}.

\begin{quote}
\emph{zhèxiē píngzi yīnggāi huíshōu yīnwèi hái néng yòng} — These bottles should be recycled, because they can still be used.
\end{quote}

Its characters come later; for now it is a word for the ear.

\begin{cousinweb}
\emph{huíshōu} is built from \zh{回} 'to turn back' and \emph{shōu} 'to gather', both inherited from Middle Chinese, \emph{hwoj} and \emph{syuw} (Guangyun, in Baxter's transcription), which glosses the second 'to collect'. Recycling is a modern sense.
\end{cousinweb}

\begin{grammarlens}[title={What you've built}]
One more word for giving reasons.
\end{grammarlens}

\subsection*{Your turn}
\begin{itemize}
\raggedright
  \item \emph{Say it:} \emph{huíshōu}
  \item \emph{Say it:} \emph{huíshōu}, once more
  \item \emph{Say it:} \emph{zhèxiē píngzi yīnggāi huíshōu yīnwèi hái néng yòng} — These bottles should be recycled, because they can still be used.
  \item \emph{Recall:} say the Mandarin for to sort into kinds, then the Mandarin for the whole family, then say \emph{huíshōu} again
\end{itemize}

\begin{tcolorbox}[breakable,colback=teal!4,colframe=teal!35!black,title={Before you move on}]
What does huíshōu mean? (\textquotedblleft{}To recycle\textquotedblright{}.) Say it once more.
\end{tcolorbox}
\section[\texorpdfstring{\zh{全国} (quánguó)}{全国 (quánguó)}]{\zh{全国} (quánguó) — the whole country; nationwide}

\label{lesson:ZH-C381-quanguo}



\begin{quote}
Before the new one: say the Mandarin for the whole family, then the Mandarin for to recycle.
\end{quote}

\subsection*{You'll want to know: \zh{全国}}
\textbf{\zh{全国}} — \emph{quánguó} — \textquotedblleft{}the whole country; nationwide\textquotedblright{}.

\begin{quote}
\textbf{\zh{今天全国都很冷}} \emph{jīntiān quánguó dōu hěn lěng} — It's cold all over the country today.
\end{quote}

\begin{cousinweb}
\textbf{\zh{全国}} (\emph{quánguó}) is built from \zh{全} 'whole' and \zh{国} 'a country', both inherited from Middle Chinese, \emph{dzjwen} and \emph{kwok} (Guangyun, in Baxter's transcription), which glosses the second 'a state'.
\end{cousinweb}

\begin{grammarlens}[title={What you've built}]
One more word for giving reasons.
\end{grammarlens}

\subsection*{Your turn}
\begin{itemize}
\raggedright
  \item \emph{Say it:} \emph{quánguó}
  \item \emph{Say it:} \emph{quánguó}, once more
  \item \emph{Say it:} \emph{jīntiān quánguó dōu hěn lěng} — It's cold all over the country today.
  \item \emph{Recall:} say the Mandarin for the whole family, then the Mandarin for to recycle, then say \emph{quánguó} again
\end{itemize}

\begin{tcolorbox}[breakable,colback=teal!4,colframe=teal!35!black,title={Before you move on}]
What does \zh{全国} mean? (\textquotedblleft{}The whole country; nationwide\textquotedblright{}.) Say it once more.
\end{tcolorbox}
\section[\texorpdfstring{\zh{安全感} (ānquángǎn)}{安全感 (ānquángǎn)}]{\zh{安全感} (ānquángǎn) — a sense of security}

\label{lesson:ZH-C381-anquangan}



\begin{quote}
Before the new one: say the Mandarin for to recycle, then the Mandarin for nationwide.
\end{quote}

\subsection*{You'll want to know: \zh{安全感}}
\textbf{\zh{安全感}} — \emph{ānquángǎn} — \textquotedblleft{}a sense of security\textquotedblright{}.

\begin{quote}
\textbf{\zh{在这个小区里我很有安全感}} \emph{zài zhège xiǎoqū lǐ wǒ hěn yǒu ānquángǎn} — In this estate I feel really safe.
\end{quote}

\begin{cousinweb}
\textbf{\zh{安全感}} (\emph{ānquángǎn}) is built from \zh{安全} 'safe' and \zh{感} 'to feel', all inherited from Middle Chinese, \emph{'an}, \emph{dzjwen} and \emph{komX} (Guangyun, in Baxter's transcription), which glosses the last 'to move'.
\end{cousinweb}

\begin{grammarlens}[title={What you've built}]
One more word for giving reasons.
\end{grammarlens}

\subsection*{Your turn}
\begin{itemize}
\raggedright
  \item \emph{Say it:} \emph{ānquángǎn}
  \item \emph{Say it:} \emph{ānquángǎn}, once more
  \item \emph{Say it:} \emph{zài zhège xiǎoqū lǐ wǒ hěn yǒu ānquángǎn} — In this estate I feel really safe.
  \item \emph{Recall:} say the Mandarin for to recycle, then the Mandarin for nationwide, then say \emph{ānquángǎn} again
\end{itemize}

\begin{tcolorbox}[breakable,colback=teal!4,colframe=teal!35!black,title={Before you move on}]
What does \zh{安全感} mean? (\textquotedblleft{}A sense of security\textquotedblright{}.) Say it once more.
\end{tcolorbox}
